\chapter{Trace minerals}

\section{Learning objectives}
Explain the roles of iron, zinc, copper, iodine, selenium, fluoride, and other trace elements; distinguish deficiency risk from laboratory noise; and recognize the hazards of concentrated food sources and mineral stacking.

\section{Iron}
Iron assessment and replacement should follow current hematology and nutrition guidance; the NIH fact-sheet collection provides updated professional summaries. \cite{odsList}
Iron is needed for hemoglobin, myoglobin, energy metabolism, and several enzymes. Heme iron comes from meat and seafood; non-heme iron comes from beans, lentils, tofu, greens, nuts, seeds, and fortified grains. Vitamin C can improve non-heme absorption, while tea, coffee, calcium, and phytate can reduce absorption around a meal. Iron deficiency can cause anemia, impaired attention, and fatigue, but iron supplements should not be used to explain every tired feeling. Menstruating people, pregnant people, infants, frequent blood donors, and people with gastrointestinal blood loss or malabsorption are higher-risk groups. Excess iron can damage organs and is dangerous for children. \cite{odsIron}

\section{Zinc and copper}
Zinc supports immune function, wound healing, DNA and protein synthesis, and taste. Meat, seafood, dairy, beans, nuts, and fortified grains provide it. Long-term high-dose zinc can cause copper deficiency and anemia. Copper supports connective tissue, iron metabolism, nervous-system function, and energy production; seafood, nuts, seeds, legumes, and organ meats provide it. Too much copper is toxic, while deficiency can occur with malabsorption or excessive zinc.

\section{Iodine and selenium}
Iodine is essential for thyroid hormone production and fetal brain development. Iodized salt, dairy, seafood, and eggs are common sources; the iodine content of seaweed varies greatly. Both deficiency and excess can disturb thyroid function. Selenium contributes to antioxidant enzymes and thyroid metabolism. Seafood, meat, eggs, dairy, and grains provide it; Brazil nuts can contain very high amounts, so frequent large servings may cause excess.

\section{Manganese, chromium, molybdenum, and fluoride}
These minerals participate in enzyme systems, metabolism, sulfur-amino-acid processing, and tooth and bone mineralization. Deficiencies are uncommon in people eating varied diets. Fluoride is not a vitamin supplement: appropriate exposure through fluoridated water and dental products helps prevent cavities, whereas swallowing concentrated fluoride is hazardous. Chromium supplements have not become a general treatment for diabetes.

\warning{Seaweed, organ meats, Brazil nuts, and concentrated mineral products can deliver unusually high doses. Food variety is safer than relying on a single ``superfood.''}

\section{Clinical reasoning}
Iron homeostasis is controlled mainly at intestinal absorption through hepcidin, the liver-derived hormone that limits ferroportin-mediated iron export. Inflammation raises hepcidin and can produce functional iron restriction even when stores are present. Ferritin rises with inflammation, liver disease, and infection; hemoglobin, mean corpuscular volume, transferrin saturation, reticulocyte response, and clinical context may also be needed. In adults with iron-deficiency anemia, occult gastrointestinal blood loss must be considered rather than attributed automatically to diet.

Zinc deficiency may cause poor wound healing, dermatitis, alopecia, diarrhea, altered taste, and impaired growth, but these signs overlap with many conditions. Prolonged high-dose zinc can impair copper status. Copper deficiency may present with anemia, neutropenia, sensory ataxia, or myelopathy and can resemble B12 deficiency. Iodine nutrition is population- and region-dependent, and urinary iodine is useful for surveillance but not a simple individual diagnostic test. Unexpected trace-element results should be confirmed and interpreted with a clinician.
